\chapter*{Resumen}

Este Trabajo de Fin de Grado presenta el diseño y desarrollo de un videojuego 
aplicando de forma práctica los conocimientos adquiridos durante la titulación 
en Ingeniería del Software. El objetivo principal no reside únicamente en la 
creación de una demo funcional, sino en utilizar el desarrollo del videojuego 
como caso de estudio para integrar metodologías, herramientas y principios 
propios de un entorno profesional de desarrollo software.

Para ello, se desarrolla Via Inferni, un videojuego roguelike 2D implementado 
con el motor Unity e inspirado en la estructura narrativa de El Infierno de 
Dante Alighieri. A lo largo del proyecto se abordan áreas fundamentales de la 
Ingeniería del Software como la gestión de proyectos, el diseño de arquitecturas 
modulares, la gestión de la configuración y del código fuente, así como la 
planificación y organización del desarrollo.

El resultado es una prueba de concepto modular y escalable que implementa 
sistemas como generación procedural de niveles, gestión de entidades, combate 
en tiempo real y progresión por partidas (runs), demostrando la viabilidad de 
aplicar principios y metodologías de Ingeniería del Software al desarrollo de 
sistemas interactivos complejos.